\documentclass[a4paper]{article}

%use the english line for english reports
%usepackage[english]{babel}
\usepackage[portuguese]{babel}
\usepackage[utf8]{inputenc}
\usepackage{indentfirst}
\usepackage{graphicx}
\usepackage{verbatim}


\begin{document}

\setlength{\textwidth}{16cm}
\setlength{\textheight}{22cm}

\title{\Huge\textbf{Q!nto}\linebreak\linebreak\linebreak
\Large\textbf{Relatório Intercalar}\linebreak\linebreak
\linebreak\linebreak
\includegraphics[scale=0.1]{feup-logo.png}\linebreak\linebreak
\linebreak\linebreak
\Large{Mestrado Integrado em Engenharia Informática e Computação} \linebreak\linebreak
\Large{Programação em Lógica}\linebreak
}

\author{
João Miguel Fidalgo Esteves Nogueira - up201303882 \\
Luís Miguel da Costa Oliveira - up201304515 \\
\linebreak\linebreak \\
 \\ Faculdade de Engenharia da Universidade do Porto \\ Rua Roberto Frias, s\/n, 4200-465 Porto, Portugal \linebreak\linebreak\linebreak
\linebreak\linebreak\vspace{1cm}}

\maketitle
\thispagestyle{empty}

%************************************************************************************************
%************************************************************************************************

\newpage

%Todas as figuras devem ser referidas no texto. %\ref{fig:codigoFigura}
%
%%Exemplo de código para inserção de figuras
%%\begin{figure}[h!]
%%\begin{center}
%%escolher entre uma das seguintes três linhas:
%%\includegraphics[height=20cm,width=15cm]{path relativo da imagem}
%%\includegraphics[scale=0.5]{path relativo da imagem}
%%\includegraphics{path relativo da imagem}
%%\caption{legenda da figura}
%%\label{fig:codigoFigura}
%%\end{center}
%%\end{figure}
%
%
%\textit{Para escrever em itálico}
%\textbf{Para escrever em negrito}
%Para escrever em letra normal
%``Para escrever texto entre aspas''
%
%Para fazer parágrafo, deixar uma linha em branco.
%
%Como fazer bullet points:
%\begin{itemize}
	%\item Item1
	%\item Item2
%\end{itemize}
%
%Como enumerar itens:
%\begin{enumerate}
	%\item Item 1
	%\item Item 2
%\end{enumerate}
%
%\begin{quote}``Isto é uma citação''\end{quote}


%%%%%%%%%%%%%%%%%%%%%%%%%%
\section{Descrição detalhada do Jogo}
\subsection{História}

“Q!nto” é um jogo de tabuleiro recente, criado no ano de 2014 no qual podem participar entre 2 e 4 jogadores. Foi criado por Gene Mackles e lançado pela PDG games

Este jogo faz parte de uma coleção de três jogos chamada Triple Play que é uma coleção de jogos de tabuleiro.

\subsection{Conteúdo do Jogo}

Estão incluídas no Jogo 60 cartas divididas como mostra a figura à direita.

\subsection{Regras de Jogo}

Existem três versões para jogar o Q!into, sendo elas o
\textit{Classic},
\textit{Plus} e
\textit{Light}. Neste projeto usaremos apenas a versão Light que passamos a explicar a seguir:

Inicialmente cada jogador retira uma carta do baralho para escolher quem começa a jogar. O jogador cuja carta tem o naipe mais alto começa, sendo que a hierarquia é a definida na imagem acima (de cima para baixo). O jogo começa dividindo igualmente o baralho pelos jogadores presentes.

As cartas que cada jogador recebeu serão a pilha de cartas com que ele vai poder jogar. Estas ficarão viradas para baixo e cada jogador retira da pilha de cartas que lhe foi atribuída as primeiras 5. 

\textbf{Noção base:} Uma linha são 2, 3, 4 ou 5 cartas numa linha ou coluna na qual as cartas têm a mesma côr ou cores diferentes e na qual as cartas têm a mesma forma ou formas diferentes. Uma linha com 5 cartas é um Q!nto.
O primeiro jogador começa o jogo. A seguir joga o jogador à sua esquerda. Quando na vez de um jogador, este pode fazer uma de três coisas:

\begin{enumerate}
\item 		Adicionar 1, 2, 3, 4 ou 5 peças à grelha numa única linha reta e depois retirar do topo da sua pilha de cartas o número necesário para voltar a ficar com 5 na mão*.
Se, numa das jogadas, o jogador completar um ou mais Q!ntos, este deve voltar a completar a sua mão* e deve jogar outra vez.
\item		Trocar uma ou mais cartas da sua mão por wildcards já jogadas e/ou a carta Q!nto de forma a que as cartas trocadas continuem a funcionar naquela posição.
\item		Passar a vez... e trocar algumas, todas ou nenhuma das susas cartas colocando-as no fundo da sua pilha e retirando do topo o número de cartas necessário para voltar a ficar com 5*.
\end{enumerate}

*ou menos, se não houver cartas suficientes para perfazer 5.
O jogo termina quando houver 1 jogador que fica sem cartas, sendo este o vencedor.
Nota: Uma carta “wild” ou “Q!nto” pode representar cartas diferentes em diferentes direções. No caso de representar coisas diferentes e não ter na mão uma carta que também as consiga representar, então, essa carta é insubstituível.


%%%%%%%%%%%%%%%%%%%%%%%%%%
\section{Visualização do Tabuleiro}

Descrever a forma de visualização do tabuleiro em modo de texto e o(s) predicado(s) Prolog construídos para o efeito.
Deve ser incluída pelo menos uma imagem correspondente ao output produzido pelo predicado de visualização.


%%%%%%%%%%%%%%%%%%%%%%%%%%
\section{Movimentos}

Elencar os movimentos (tipos de jogadas) possíveis e definir os cabeçalhos dos predicados que serão utilizados (ainda não precisam de estar implementados).


\end{document}
